\documentclass[conference]{IEEEtran}

\usepackage[T1]{fontenc}
\usepackage[utf8]{inputenc}
\usepackage{cite}
\usepackage{graphicx}
\usepackage{booktabs}
\usepackage{multirow}
\usepackage{array}
\usepackage{xurl}
\usepackage[hidelinks]{hyperref}
\Urlmuskip=0mu plus 1mu

\title{Hardware Security}

\author{\IEEEauthorblockN{Author}
\IEEEauthorblockA{Hardware Security Survey}}

\begin{document}

\maketitle

\begin{abstract}
This survey reviews defense mechanisms for Arm and RISC-V confidential
computing, confidential-computing network and I/O data paths, and ISA-level
hardware design defenses. The central observation is that modern hardware
security cannot be evaluated by CPU isolation alone. A confidential workload is
protected only when execution domains, physical-memory ownership, device
assignment, DMA translation, interrupt routing, link protection, and attestation
evidence compose into a coherent platform boundary. The survey therefore
connects Arm CCA, RISC-V CoVE/AP-TEE, trusted I/O protocols, accelerator TEEs,
and hardware mechanisms such as MTE, PMP, capability machines, and CFI
extensions under a common systems-oriented taxonomy.
\end{abstract}

\begin{IEEEkeywords}
hardware security, confidential computing, trusted execution environments,
Arm CCA, RISC-V, attestation, I/O protection
\end{IEEEkeywords}

\section{Introduction}

Hardware security has shifted from isolated processor features toward
platform-level protection contracts. Modern systems must coordinate execution
domains, memory ownership, device access, interrupt delivery, and attestation
evidence across the whole machine. This survey organizes the current literature
around that composition problem, using Arm and RISC-V as the main architectural
lines.

The scope is deliberately narrower than a general hardware-security SoK. We
focus on three defense-oriented lanes. The first lane covers confidential
computing mechanisms in Arm and RISC-V systems, including Arm TrustZone and CCA
as well as RISC-V enclave lineage, CoVE, AP-TEE, and TVM lifecycle management.
The second lane covers network, I/O, and data-path defense in
confidential-computing deployments, including SPDM, TDISP, PCIe IDE,
IOMMU/IOPMP/SMMU mediation, trusted MSI, SmartNIC/DPU-style offload, and
accelerator TEEs. The third lane covers ISA and hardware design defenses that
provide defense-in-depth for those platforms: memory tagging, page-table
permissions, PMP/ePMP, capability/tagging architectures, hardware CFI,
memory-encryption taxonomy, and debug or trace controls.

This scope also defines what the survey does not attempt to cover in the
current version. Side channels, physical leakage, fault injection, Rowhammer,
and attack-only vulnerability catalogs are important hardware-security topics,
but they are not treated as independent research lanes here. They are mentioned
only when they affect the threat-model boundary of a confidential-computing
mechanism. Similarly, ordinary network-security mechanisms such as firewalls,
generic IDS/IPS systems, routing security, TLS, or VPNs are excluded unless they
directly protect a confidential workload's trusted network endpoint, device
assignment path, or secure offload boundary.

The paper is organized as follows. The Arm section explains the transition from
TrustZone-style TEEs to CCA and Realm-based protection. The RISC-V section
connects earlier enclave systems to CoVE/AP-TEE confidential VMs. The
confidential I/O section then treats devices, fabrics, and network data paths
as first-class parts of the confidential-computing boundary. Finally, the
hardware-design section surveys cross-architecture ISA and microarchitectural
defenses that complement those confidential-computing mechanisms.

\section{Methodology and Evidence Ranking}

This survey separates evidence strength from topical relevance. A source may
be important to the taxonomy while still being unsuitable for a mechanism-level
claim. We therefore use the following evidence classes when deciding how a
reference may support the text:

\begin{table}[htbp]
\centering
\caption{Evidence Classes and Citation Rules}
\label{tab:evidence_ranking}
\begin{tabular}{@{}p{0.12\linewidth}p{0.36\linewidth}p{0.42\linewidth}@{}}
\toprule
\textbf{Class} & \textbf{Evidence type} & \textbf{Allowed use in this survey} \\ \midrule
E0 & Official ratified specification, architecture manual, RFC, or standards document & Normative mechanism and status claims, with no experimental claims unless the source provides them. \\
E1 & Peer-reviewed primary research paper & Mechanism, evaluation, and threat-model claims supported by the reviewed paper and its recorded source/PDF status. \\
E2 & Survey, SoK, or secondary review & Taxonomy, terminology, and coverage guidance; original mechanism claims are traced back to primary papers or specifications. \\
E3 & Public draft, release candidate, or not-ratified specification & Current standardization direction, always labelled as draft, release candidate, or not ratified. \\
E4 & Vendor, industry, or product documentation & Industry evidence for product behavior or deployment practice, not proof of general security unless independently supported. \\
E5 & Metadata-only, gated, blocked-PDF, or HTML-only evidence & Existence, source status, or backlog classification only; no substantive mechanism claim unless promoted or paired with stronger public evidence. \\ \bottomrule
\end{tabular}
\end{table}

The citation rules are intentionally conservative. First, all mechanism claims
prefer E0 or E1 sources. Second, E2 sources such as surveys and SoKs are used as
maps of the literature, not as substitutes for original mechanisms. Third, E3
drafts are useful for current architecture direction, but the text must state
their draft or not-ratified status. Fourth, E4 industry documents are treated
as deployment or product evidence and kept separate from peer-reviewed system
claims. Fifth, E5 sources, including metadata-only \texttt{rw...} candidates and
member-gated documents, do not support substantive text claims unless a public
source supports the same claim. When a source has mixed properties, such as a
public metadata page for a gated specification or an HTML-only survey, the more
restrictive citation rule controls the claim.

\section{Arm Confidential Computing}

Confidential computing in the Arm ecosystem addresses a stronger threat model than classical trusted-execution environments. The central problem is not only how to isolate a secure software stack from ordinary software, but how to protect the confidentiality of a workload even when privileged software still manages scheduling, memory allocation, and devices~\cite{arm_trustzone_whitepaper,pinto2019trustzone,li2022cca,arm_cca_spec,arm_rme_spec}. In Arm platforms, this evolution is visible in the transition from TrustZone-based TEEs to the Arm Confidential Compute Architecture (CCA), which extends the protection boundary into physical-memory ownership, device mediation, and attestation. The following sections survey this transition, covering TrustZone-based systems, Realm-based isolation, hardware enforcement in the memory and I/O path, and the attestation mechanisms needed to make protected execution usable in practice.

\subsection{Threat Model and Design Objectives}

\subsubsection{From Trusted Services to Host-Resilient Confidentiality}

The distinction between TrustZone-era TEEs and Arm CCA is fundamentally a distinction in threat model. In a classical TrustZone deployment, the secure-world firmware and runtime environment are themselves part of the trusted computing base, and the main goal is to isolate security services from the normal world~\cite{arm_trustzone_whitepaper,pinto2019trustzone}. In the confidential-computing setting addressed by Arm CCA, the protected workload must remain confidential even though a host stack continues to manage cores, memory allocation, and devices~\cite{li2022cca,arm_cca_spec,arm_rme_spec}. The architecture therefore aims to separate resource management from data visibility more explicitly than conventional TEE models do.

\subsubsection{Core Design Dimensions}

The cited Arm CCA documents and systems papers show that this protection problem decomposes into several recurring design dimensions. The first is execution-domain structure: the architecture needs a protected execution context that is distinct from both ordinary software and the management software responsible for scheduling and provisioning~\cite{li2022cca,arm_cca_spec,arm_rme_spec}. The second is memory ownership: physical memory must carry hardware-enforced ownership metadata rather than depending only on hypervisor-maintained bookkeeping~\cite{arm_cca_spec,arm_rme_spec,linux_arm_cca_doc}. The third is system integration: I/O, accelerators, interrupts, and attestation all have to reflect the same protection boundary if confidentiality is to survive contact with the whole machine~\cite{arm_smmu_spec,acai2023,bertschi2026devlore,arm_rmm_spec,eat_rfc}. These dimensions provide a useful organizing lens for the remainder of the chapter.

\subsection{From TrustZone-Based TEEs to Arm CCA}

\subsubsection{TrustZone as the Classical Arm Isolation Model}

\textbf{TrustZone} partitions an Arm system into Secure and Non-secure worlds, with the Non-secure bit propagated through the SoC interconnect and memory system~\cite{arm_trustzone_whitepaper,pinto2019trustzone}. This model has been widely used to build TEEs for mobile, embedded, and IoT devices, where a relatively small secure-world software stack mediates access to protected resources. Representative systems such as \textbf{TrustShadow} show that TrustZone can protect unmodified applications from an untrusted normal-world operating system by relocating sensitive execution into the secure world~\cite{guan2017trustshadow}. At the same time, survey and SoK work shows that TrustZone-assisted TEEs remain exposed to recurring weaknesses in interfaces, implementation choices, and trusted computing base growth~\cite{pinto2019trustzone,cerdeira2020trustzone}.

For hardware-security surveys, the importance of TrustZone is twofold. First, it demonstrates that Arm's security story has long depended on system-wide enforcement rather than on CPU privilege levels alone. Second, it shows the limits of a two-world model once the secure world accumulates more functionality: as more services migrate into the trusted side, the TCB grows and the isolation boundary becomes harder to reason about.

\subsubsection{Arm CCA and the Realm Model}

\textbf{Arm CCA} extends this design through the \textbf{Realm Management Extension (RME)}, introducing a model in which confidential workloads execute inside \textbf{Realms} that are intended to remain opaque to the host software stack~\cite{li2022cca,arm_cca_spec,arm_rme_spec}. The associated \textbf{Realm Management Monitor (RMM)} manages Realm creation, execution, and teardown~\cite{arm_rmm_spec,wu2024rmm}. Compared with TrustZone, the key change is not merely the addition of another execution environment; it is the explicit separation between resource management and data ownership. The host still provisions CPUs, memory, and devices, but it is not meant to inspect Realm memory or register state as part of normal management.

This change has already motivated systems work above the raw architecture. \textbf{SHELTER} uses Arm CCA to extend isolation into user space~\cite{zhang2023shelter}. \textbf{RContainer} treats CCA hardware primitives as a basis for container isolation~\cite{zhou2025rcontainer}. \textbf{virtCCA} studies how CCA-style protection can be virtualized in conjunction with existing Arm trust mechanisms~\cite{xu2026virtcca}. Taken together, these works show that Arm CCA is being treated not only as an architectural feature set, but also as a substrate for higher-level confidential-computing software stacks.

\begin{table*}[t]
\centering
\caption{Arm CCA Mechanism Map}
\label{tab:arm_cca_mechanism_map}
\begin{tabular}{@{}p{0.18\linewidth}p{0.28\linewidth}p{0.24\linewidth}p{0.20\linewidth}@{}}
\toprule
\textbf{Mechanism} & \textbf{Role in CCA} & \textbf{Security Property} & \textbf{Survey Boundary} \\ \midrule
Realm security state & Separates Realm execution from Non-secure, Secure, and Root worlds & Host cannot directly inspect Realm execution state under the CCA model & Confidential VM / Realm isolation, not a generic process sandbox \\
RMM & Mediates Realm creation, entry, exit, destruction, and memory lifecycle operations & Keeps host resource management outside the Realm data boundary & RMM is in the CCA TCB; correctness matters as much as hardware support \\
GPT/GPC & Records and checks physical granule ownership and security state & Blocks accesses inconsistent with the assigned physical-address space & Access-control and ownership metadata, not by itself memory encryption \\
PAS and RIPAS & Represent physical-address-space assignment and Realm IPA state & Makes delegation, acceptance, destruction, and sharing explicit state transitions & Needed to reason about stale mappings, sharing, and reclaim races \\
RMI & Host-to-RMM management interface for Realm and granule lifecycle & Lets an untrusted host manage resources through constrained operations & Management authority is preserved but data visibility is removed \\
RSI & Realm service interface used by Realm software to request CCA services & Gives Realm software a controlled way to interact with the trusted CCA layer & Must be distinguished from host-facing RMI calls \\
Attestation evidence & Reports platform, RMM, and initial Realm state to a verifier & Enables key release and remote trust decisions & Evidence usefulness depends on verifier policy and measured claims \\
SMMU and device mediation & Constrains DMA and device-visible memory for Realm-related I/O & Prevents CPU-only isolation from being bypassed by devices & Requires confidential I/O mechanisms beyond base Realm memory ownership \\ \bottomrule
\end{tabular}
\end{table*}

Table~\ref{tab:arm_cca_mechanism_map} fixes the mechanism vocabulary used in the rest of this survey. The main distinction is between \emph{management authority} and \emph{data authority}. CCA still allows the host to allocate CPUs, schedule execution, and request lifecycle transitions through RMI, but GPT/GPC checks, RMM state, PAS/RIPAS transitions, and Realm-facing RSI calls constrain what that management software can observe or mutate~\cite{arm_cca_spec,arm_rme_spec,arm_rmm_spec,linux_arm_cca_doc}. This table also prevents a common classification error: CCA memory ownership metadata should not be described as memory encryption, and RMM lifecycle logic should not be collapsed into ordinary page-table policy.

\subsection{Hardware Enforcement in the Memory and I/O Path}

\subsubsection{Granule Ownership and Physical-Memory Checks}

The stronger isolation model of Arm CCA depends on hardware checks in the physical-memory path. The \textbf{Granule Protection Table (GPT)} records the ownership and security state of memory granules, while the \textbf{Granule Protection Check (GPC)} validates physical accesses against that metadata~\cite{arm_cca_spec,arm_rme_spec,linux_arm_cca_doc}. This is a significant departure from designs in which isolation is represented only by software-maintained page tables or a small number of coarse protected regions. In the CCA model, physical-memory ownership becomes a hardware-checked property that constrains even privileged software.

This ownership model is also what makes CCA suitable for dynamic workloads. Confidential workloads may be created, destroyed, resized, or migrated, and their protected memory footprint may change over time. Recent work such as \textbf{SHELTER}, \textbf{NanoZone}, and \textit{More Granular, Less Trust} exploits this finer-grained ownership model to study user-space and intra-process isolation under an untrusted manager~\cite{zhang2023shelter,liu2025nanozone,liu2025lesstrust}. The common theme across these systems is that the architecture is no longer treating protected memory as a single static partition; instead, it supports repeated ownership transitions while preserving hardware-enforced separation.

\subsubsection{Lifecycle Management and Ownership Transitions}

The granule model is closely tied to lifecycle management. The \textbf{RMM} and the Arm CCA software model treat Realm creation, memory delegation, execution, and teardown as explicit state transitions rather than as ad hoc software conventions~\cite{arm_rmm_spec,linux_arm_cca_doc}. This matters because confidential-computing platforms must repeatedly move memory between protection states while preserving consistency between CPU execution, physical-memory ownership, and the host's remaining management authority. The design and verification work of Li et al.~\cite{li2022cca} and later verification-oriented work on the \textbf{RMM}~\cite{wu2024rmm} both reinforce the point that lifecycle state management is part of the security argument, not merely an implementation detail.

\subsubsection{Securing Devices, Accelerators, and Memory Fabrics}

CPU-side isolation is insufficient if \textbf{DMA}-capable devices can still reach protected memory directly. Arm therefore extends the confidentiality story into the I/O path through \textbf{SMMUv3}-based mediation and realm-aware device assignment~\cite{arm_smmu_spec,acai2023,bertschi2026devlore}. This matters because many confidential workloads depend on accelerators, storage engines, network adapters, or packet-processing units. If device access is not mediated under the same protection regime as CPU access, the device path becomes the obvious route around the confidentiality boundary.

Recent systems research makes this point concrete. \textbf{ACAI} studies how accelerator execution can be protected using Arm CCA mechanisms~\cite{acai2023}. \textbf{CAGE} extends this line by adapting GPU and FPGA accelerator workflows to Realm-style execution with shadow tasks and GPC/GPT-backed checks~\cite{wang2026cage}. \textbf{PORTAL} addresses a neighboring device-access problem for mobile SoCs: it uses CCA memory isolation to create a protected plaintext device-access region instead of encrypting every transfer through an untrusted shared buffer~\cite{sang2025portal}. \textbf{Devlore} focuses on device-interrupt protection for confidential VMs, highlighting that interrupt delivery and device ownership are part of the same protection problem~\cite{bertschi2026devlore}. Earlier Arm GPU TEE work such as \textbf{StrongBox} is useful as historical context because it shows how TrustZone and Stage-2 translation were used to protect unified-memory GPU tasks before CCA-style Realm mechanisms became the main confidential-computing target~\cite{deng2022strongbox}. Related work on \textbf{CXL}, \textbf{PCIe IDE}, heterogeneous memory descriptions, and remote paging is background substrate: it shows that disaggregated memory and high-speed interconnects affect where protected data resides and which agents can observe it, but those sources are not themselves Arm CCA mechanisms or complete trusted-I/O evidence~\cite{cxl_spec,pcie_ide,acpi_spec,gouk2022directcxl,zhong2024cxltiers,wang2025odrp}. For confidential workloads, the integrity of the protection boundary therefore depends on the fabric and device path as much as on the CPU core.

\subsection{Software Stacks and Deployment Models}

\subsubsection{User-Space, Containers, and Virtualized Realms}

One reason Arm CCA has attracted attention is that it can support several deployment models above the raw hardware interface. \textbf{SHELTER} explores user-space isolation with Arm CCA hardware primitives~\cite{zhang2023shelter}. \textbf{RContainer} extends the same substrate toward confidential containers~\cite{zhou2025rcontainer}. \textbf{virtCCA} studies virtualized CCA deployments that combine Realm-style protection with existing Arm trust infrastructure~\cite{xu2026virtcca}. \textbf{Aster} brings the same question into the Android ecosystem by mapping Android Virtualization Framework protected VMs to CCA Realms and identifying the boot, attestation, rollback, and memory-protection changes needed for mobile protected workloads~\cite{kuhne2026aster}. These works are useful in a survey because they show that CCA is not restricted to one software packaging model; rather, it is being used as a common basis for protected processes, containers, virtualized workloads, and mobile protected VMs.

\begin{table*}[t]
\centering
\caption{Arm CCA Deployment Taxonomy}
\label{tab:arm_cca_deployment_taxonomy}
\scriptsize
\begin{tabular}{@{}p{0.14\linewidth}p{0.15\linewidth}p{0.15\linewidth}p{0.18\linewidth}p{0.12\linewidth}p{0.12\linewidth}@{}}
\toprule
\textbf{Category} & \textbf{Representative work} & \textbf{Protected object} & \textbf{CCA mechanism used} & \textbf{Evidence role} & \textbf{Main limitation} \\ \midrule
User-space isolation & SHELTER & Process or application compartment & Realm/GPT/GPC-backed user-space isolation & E1 primary paper & Not a general VM or container deployment. \\
Container isolation & RContainer & Container workload & CCA primitives beneath container packaging & E1 primary paper & Deployment maturity and orchestration integration remain open. \\
Virtualized CCA & virtCCA & Virtualized Realm/CCA stack & CCA combined with existing Arm trust infrastructure & E3 arXiv/preprint evidence & Nested trust and virtualization complexity. \\
Mobile protected VM & Aster & Android AVF protected VM & Realm-world pVM execution with DICE, policy, rollback, and memory-protection integration & E1 emerging mobile CCA evidence & Evaluated with emulator/Armv8.2 approximation rather than native CCA Android hardware. \\
Intra-process isolation & NanoZone; More Granular, Less Trust & Fine-grained compartment inside a process & Granule ownership and Realm-style memory protection & E3 plus E1 & Granularity and draft/preprint status vary by work. \\
CVM maintenance & CPC & Migration, snapshot, and resource reclamation procedure & Guest-side confidential procedure plus CPTI-style intra-CVM isolation & E1 primary paper & Arm CCA evidence uses official simulator rather than production hardware. \\
Inter-CVM sharing & CAEC & Communication between confidential VMs & Attestation plus controlled sharing over CCA substrate & E3 arXiv evidence & Not an Arm official interface; sharing policy remains research. \\
Device and accelerator access & PORTAL; CAGE; Devlore & SoC device, accelerator workflow, or interrupt path & SMMU, GPT/GPC, GPC-backed checks, and Realm-aware mediation & E1 plus E3 & Device identity and complete trusted-I/O lifecycle are still incomplete. \\ \bottomrule
\end{tabular}
\end{table*}

Table~\ref{tab:arm_cca_deployment_taxonomy} separates these systems by the
software object they protect rather than by paper name. This matters because a
process compartment, container, nested CCA stack, inter-CVM channel, and device
path stress different parts of the architecture. The table also keeps evidence
roles visible: peer-reviewed systems can support their own mechanisms and
measurements, while arXiv or pre-standardization work is treated as emerging
research evidence rather than as proof of mature deployment.

\textbf{CPC} adds a separate lifecycle dimension to this deployment taxonomy:
CVM maintenance. Instead of giving an untrusted host direct access to guest
private state or moving maintenance modules into firmware, CPC lets the host
trigger tenant-trusted maintenance procedures inside the CVM through dedicated
vCPU scheduling semantics~\cite{chen2024cpc}. This makes live migration,
snapshotting, and resource reclamation part of the confidential-computing
management problem. The paper's AMD SEV-SNP prototype provides real-machine
performance evidence, while the Arm CCA prototype is simulator-based; therefore
this survey uses CPC as peer-reviewed evidence for the maintenance/lifecycle
gap, not as an Arm official interface or production CCA migration standard.

\textbf{Aster} is treated with a similar boundary discipline. It is useful
peer-reviewed evidence that Arm CCA can be studied as an Android AVF backend,
including pVM launch policy, DICE-compatible attestation, rollback protection,
and per-pVM memory protection~\cite{kuhne2026aster,tcg_dice_2018}. It is not an Arm official
mobile profile, and its evaluation is constrained by the absence of native
CCA-enabled Android hardware. The survey therefore uses it as emerging mobile
deployment evidence rather than as production Android CCA availability.

\subsubsection{Fine-Grained Isolation Above the Base Architecture}

Later systems work also shows that Arm CCA is being pushed toward finer-grained protection goals than those envisioned by a simple confidential-VM model. \textbf{NanoZone} focuses on scalable and efficient memory protection for Arm CCA~\cite{liu2025nanozone}. \textit{More Granular, Less Trust} studies intra-process isolation under an untrusted management environment~\cite{liu2025lesstrust}. Together with \textbf{SHELTER}, these works indicate that once granule ownership and Realm mechanisms exist, researchers naturally try to repurpose them for user-space compartmentalization and finer-grained software isolation. This is an important trend because it suggests that Arm CCA is evolving from a narrowly server-oriented confidential-computing mechanism into a more general substrate for trusted isolation.

\subsection{Attestation and Trust Establishment}

\subsubsection{Attestation for Realms}

Isolation alone does not let a remote party decide whether to provision secrets to a workload. Arm CCA therefore couples protected execution with attestation evidence describing the platform state, the management monitor, and the initial Realm state~\cite{arm_cca_spec,arm_rmm_spec}. The \textbf{RATS} architecture defines the attester, verifier, relying-party, evidence, and appraisal roles, while the \textbf{Entity Attestation Token (EAT)} standard provides a general token format for carrying attestation claims in a machine-consumable form~\cite{rats_rfc,eat_rfc}. \textbf{DICE} supplies lower-level boot identity vocabulary by deriving a Compound Device Identifier from a Unique Device Secret and first mutable code measurement before mutable code runs~\cite{tcg_dice_2018}. More broadly, studies of attestation mechanisms for TEEs emphasize that the practical value of attestation depends on what is measured, how claims are encoded, and how a verifier evaluates them~\cite{menetrey2022attestation}. In this sense, attestation converts confidential execution from a local hardware property into a system property that can be consumed by remote tenants, orchestration systems, and key brokers.

\subsubsection{Broader Attestation Context in Arm Systems}

The same lesson appears in earlier TrustZone-oriented literature. Work on secure boot, trusted boot, mutual remote attestation, and runtime integrity measurement shows that deployable trust depends on evidence spanning both launch time and runtime~\cite{ling2021trustzoneatt,chen2024mraima,mao2025pdrima}. For edge and IoT platforms, \textbf{PSA Certified} provides a broader framework for hardware-rooted assurance and attestation-oriented system design~\cite{psa_certified}. Earlier attestation systems such as \textbf{SWATT}, \textbf{SMART}, and \textbf{SEDA} reinforce the same point from a wider systems perspective: evidence must be attributable to a concrete execution context and fresh enough to support a real trust decision~\cite{seshadri2004swatt,defrawy2012smart,asokan2015seda}. Within Arm CCA, attestation is therefore not an optional add-on; it is the mechanism that makes Realm isolation externally verifiable.

\subsection{Comparative Perspective}

\subsubsection{TrustZone versus Arm CCA}

The transition from \textbf{TrustZone} to \textbf{Arm CCA} is best understood as a change in both threat model and enforcement granularity. TrustZone is effective when a secure-world software stack is itself part of the trusted boundary, but Arm CCA is designed for settings in which a powerful host stack must remain outside the confidentiality boundary~\cite{arm_trustzone_whitepaper,pinto2019trustzone,li2022cca,arm_cca_spec,arm_rme_spec,arm_rmm_spec}. As a result, protection moves from a mostly binary world partition to an ownership model backed by GPT/GPC checks and Realm lifecycle management. The comparison is not that one mechanism replaces the other in all settings; rather, they target different deployment assumptions.

\begin{table}[htbp]
\centering
\caption{Architectural Comparison: TrustZone and Arm CCA}
\label{tab:tz_vs_cca}
\begin{tabular}{@{}p{0.22\linewidth}p{0.33\linewidth}p{0.35\linewidth}@{}}
\toprule
\textbf{Property} & \textbf{TrustZone} & \textbf{Arm CCA} \\ \midrule
Primary Isolation Model & Secure / Non-secure world split & Realm-based confidential execution with Root/Secure/Realm/Non-secure separation \\
Protection Granularity & Platform-defined secure-world partitioning & Granule ownership tracked through GPT/GPC \\
Trusted Software Context & Secure-world firmware and services & RMM-mediated Realm lifecycle with host outside the data boundary \\
Memory Enforcement Style & World-aware interconnect and secure memory regions & Hardware-checked ownership transitions for protected granules \\
Device Model & Typically static secure-world device handling & Realm-aware device and I/O mediation \\
Attestation Practice & Often platform- or vendor-specific & Structured evidence for Realm launch and platform state \\ \bottomrule
\end{tabular}
\end{table}

\subsubsection{Arm CCA in Comparative Context}

\textbf{AMD SEV-SNP} addresses a similar high-level problem by protecting guest workloads from a host with substantial management authority~\cite{amd_sev,amd_sev_snp}. The useful comparison is therefore not whether both systems support confidential workloads, but how they represent and enforce the protection boundary. In the Arm case, the cited architecture documents and systems papers emphasize explicit security states, Realm lifecycle management, and hardware ownership checks over physical memory~\cite{li2022cca,arm_cca_spec,arm_rme_spec}. By contrast, the comparison with \textbf{RISC-V PMP} highlights a different design axis: CCA uses table-backed ownership metadata, whereas PMP uses machine-mode configuration registers to define protected physical regions~\cite{arm_cca_spec,arm_rme_spec,riscv_privileged}. The latter remains an important primitive for smaller systems, but the former is better aligned with large, dynamic workloads that repeatedly change memory ownership and depend on coordinated device mediation.

\subsection{Design Tradeoffs and Open Challenges}

\subsubsection{Scalability, Complexity, and System Composition}

The literature surveyed above suggests that Arm CCA's strengths also imply new engineering burdens. Table-backed ownership tracking and explicit lifecycle transitions are more expressive than coarse partitioning, but they require careful coordination among the architecture, the monitor, the host software stack, and devices~\cite{li2022cca,arm_cca_spec,arm_rme_spec,arm_rmm_spec}. The emergence of systems such as \textbf{NanoZone}, \textbf{RContainer}, and \textbf{virtCCA} indicates that researchers are already exploring how to scale these mechanisms to user space, containers, and virtualized deployments without losing the clarity of the protection boundary~\cite{liu2025nanozone,zhou2025rcontainer,xu2026virtcca}. This suggests that scalability is not only a performance question; it is also a question of whether the management model remains understandable as deployments become more complex.

\subsubsection{Data-in-Use versus Data on the System Fabric}

A second recurring issue is that confidential execution cannot be reduced to CPU-local isolation. In Arm CCA specifically, this concern appears through the need to align Realm isolation with device mediation, interconnect protection, and attestation rather than through one standalone mechanism~\cite{arm_cca_spec,arm_rme_spec,arm_smmu_spec,cxl_spec,pcie_ide}. CXL and PCIe IDE are cited here as data-path and link-protection context, not as proof that Arm CCA alone secures every fabric endpoint. The architecture therefore addresses confidential computing as a composition problem: protected execution, memory ownership, I/O mediation, and trust evidence have to remain consistent across the entire platform.

\subsection{Implications for Network Devices}

The mechanisms surveyed in this section are directly relevant to modern network hardware. Network appliances increasingly combine general-purpose cores, \textbf{SmartNICs}, \textbf{DPUs}, storage engines, and accelerators on the same platform, and many of these components are \textbf{DMA}-capable. In such systems, protecting a confidential workload is not just a matter of isolating CPU execution; it also requires consistent mediation of memory ownership, device assignment, interrupts, and attestation evidence. Arm CCA is therefore particularly relevant to network equipment that performs virtualization, confidential packet processing, or multi-tenant offload, because the same platform must let the operator manage resources without automatically exposing tenant memory or accelerator state. Current accelerator and DPU work shows four complementary routes: protect the accelerator through a CCA-aware path, as in ACAI and CAGE; protect SoC devices through CCA-managed access regions, as in PORTAL; protect rack- or cloud-scale heterogeneous devices through a security controller, as in HETEE and CloudScale; or rely on device-local roots such as DPU OP-TEE/fTPM building blocks, NIC-local roots, and accelerator-local TEEs~\cite{acai2023,wang2026cage,sang2025portal,zhu2020hetee,dhar2024cloudscale,nvidia_bluefield_operation_2025,giantsidi2025tnic,vaswani2023itx}. The DPU OP-TEE/fTPM material is vendor evidence for building blocks, not proof of a complete production confidential-networking TEE.

\subsection{Summary}

The Arm literature surveyed here shows a clear progression from TrustZone-based TEEs to confidential-computing architectures with stronger assumptions about privileged adversaries. TrustZone established a system-wide isolation model based on Secure and Non-secure worlds, but its protection boundary is relatively coarse and often tied to a growing secure-world TCB. Arm CCA extends the design by introducing Realms, GPT/GPC-backed memory ownership, RMM-managed lifecycle control, device mediation through the I/O path, and structured attestation evidence~\cite{arm_trustzone_whitepaper,pinto2019trustzone,li2022cca,arm_cca_spec,arm_rme_spec,arm_rmm_spec}. A cross-cutting observation is that confidential computing in Arm is not defined by any single primitive. Its security properties emerge from the composition of memory ownership tracking, device mediation, and externally verifiable attestation, all of which must hold simultaneously if the workload is to remain protected from a powerful host environment.


\section{RISC-V Confidential Computing}

RISC-V confidential computing is best understood as the convergence of two
historically separate lines of work. The first line is the open enclave
literature, where researchers use RISC-V's extensible ISA, privilege model, and
physical-memory protection mechanisms to build small trusted monitors and
enclave runtimes. The second line is the standardization-oriented confidential
VM effort represented by CoVE and AP-TEE, where the target is no longer only an
application enclave but a tenant VM protected from an untrusted host
environment~\cite{sahita2023cove,riscv_ap_tee_2024,boubakri2025riscvtee}.
Boubakri et al. are used here as survey/background substrate; mechanism and
status claims in this section rely on the original CoVE paper, AP-TEE draft, or
other primary sources. This distinction is important for classification:
PMP-based enclaves, tagged-memory enclaves, and CoVE-style TVMs share a
motivation, but they do not expose the same lifecycle, device, or attestation
contract.

\subsection{From Enclave Research to Confidential VMs}

\subsubsection{Early RISC-V Enclave Design}

\textbf{Sanctum} is one of the earliest and most influential examples of using
RISC-V to explore strong software isolation through minimal hardware
extensions~\cite{costan2016sanctum}. Its importance is not just that it provides
an enclave system, but that it frames open hardware as a way to make TEE
assumptions explicit. Instead of treating a proprietary platform as a fixed
black box, Sanctum exposes the processor, cache, memory, and monitor assumptions
that must hold for enclave isolation.

\textbf{Keystone} then provides a more reusable open-source RISC-V TEE
framework built around a security monitor, PMP-based memory isolation, and
runtime abstractions for enclaves~\cite{lee2020keystone}. Keystone is central to
the survey because it shows the strengths and limitations of using standard
RISC-V mechanisms for protected execution. PMP gives a simple and efficient
physical-memory filter, but a small number of PMP regions and monitor-managed
state are not the same as a cloud-scale confidential-VM memory ownership model.
This is the point at which RISC-V enclave research begins to diverge from
later CoVE/AP-TEE work.

\subsubsection{Scaling and Specializing the Enclave Boundary}

Subsequent RISC-V enclave systems explore different ways to scale, specialize,
or strengthen the basic monitor-plus-memory-protection model. \textbf{TIMBER-V}
uses tag-isolated memory to support fine-grained enclaves, which is especially
relevant for embedded and compartmentalized systems where the protected unit may
be smaller than a full VM~\cite{weiser2019timberv}. \textbf{CURE} introduces
customizable resilient enclaves, emphasizing configurable isolation and
resilience tradeoffs~\cite{bahmani2021cure}. \textbf{MI6} studies secure
enclaves in the context of a speculative out-of-order processor, making clear
that the microarchitecture matters even when the ISA-level enclave abstraction
appears simple~\cite{bourgeat2019mi6}.

\textbf{Penglai} addresses a different bottleneck: scalable memory protection
for large numbers of enclaves, motivated by cloud and serverless deployment
patterns~\cite{feng2021penglai}. \textbf{SPEAR-V} proposes low-overhead
enclave primitives for RISC-V and is useful as a later point in the lineage
because it studies practical primitive design rather than only monitor
structure~\cite{schrammel2023spearv}. \textbf{Cerberus} focuses on formal and
efficient enclave memory sharing, an issue that becomes increasingly important
once protected domains need deliberate shared-memory communication rather than
complete isolation~\cite{lee2022cerberus}. \textbf{ACE} extends the discussion
toward confidential computing for embedded RISC-V systems~\cite{ozga2025ace}.

Across these works, the recurring lesson is that a TEE is not only an execution
mode. It is a contract over memory ownership, entry and exit, metadata
integrity, sharing, interrupt behavior, and attestation. Earlier enclave
systems often solve a subset of that contract for a specific threat model.
CoVE/AP-TEE attempts to standardize a larger contract for confidential VMs.

\subsection{CoVE and AP-TEE}

\subsubsection{Threat Model and Architectural Role}

RISC-V CoVE targets a host-resilient confidential-VM model. The tenant workload
executes in a trusted virtual machine while a host stack remains responsible for
resource management, scheduling, and device orchestration~\cite{sahita2023cove}.
This is conceptually close to Arm CCA in that both designs separate management
authority from data visibility, but the RISC-V path is expressed through a
different standardization structure and a more modular ecosystem. The AP-TEE
specification defines the relevant ABI and lifecycle concepts, but it remains a
draft and should be treated as a standardization work in progress rather than a
ratified architecture~\cite{riscv_ap_tee_2024}.

The core AP-TEE objects are the \textbf{Trusted Security Monitor (TSM)}, the
\textbf{TSM driver}, \textbf{Supervisor Domains}, and the \textbf{Trusted
Virtual Machine (TVM)}~\cite{riscv_ap_tee_2024}. The host retains the ability
to allocate resources, but the TSM mediates TVM creation, memory assignment,
entry, exit, and security-sensitive transitions. This moves RISC-V beyond the
simple observation that PMP can isolate memory. A confidential VM needs a
defined lifecycle, a way to donate and reclaim pages, an attestation path, and
an interface that lets untrusted management software perform useful work without
reading protected state.

\subsubsection{Memory Lifecycle and TVM State}

Memory lifecycle is the central mechanism that differentiates CoVE/AP-TEE from
earlier PMP-only enclave explanations. The draft AP-TEE model includes memory
\textbf{donation}, \textbf{reclaim}, and \textbf{sharing} semantics, which are
needed because a TVM is dynamic: pages are created, mapped, unmapped, shared for
communication, and eventually returned to the host~\cite{riscv_ap_tee_2024}.
These transitions are security-critical. If a page can move between host and TVM
ownership without explicit state checks, stale mappings, unclean metadata, or
ambiguous sharing state can break the confidentiality boundary.

This is also where RISC-V confidential computing should not be reduced to
memory encryption. The AP-TEE draft primarily specifies ownership, lifecycle,
and access-control semantics; it does not by itself replace a memory-encryption
and replay-protection taxonomy~\cite{henson2014memory}. A complete platform may
combine access control, encryption, integrity protection, and attestation, but
those mechanisms answer different questions. Access control asks which agent
may reach a physical page. Encryption asks what an observer of raw memory can
learn. Integrity and replay protection ask whether stale or modified data can be
detected. A precise survey has to keep those layers separate.

\subsubsection{Attestation and Evidence}

AP-TEE also requires an attestation story because tenants need external
evidence before provisioning secrets to a TVM~\cite{riscv_ap_tee_2024}. The
evidence has to bind the platform, TSM, TVM initial state, and relevant
configuration to a verifier-visible claim. This mirrors the broader attestation
literature: a TEE claim is useful only if the measured state is well-defined,
fresh, tied to the execution context that will receive secrets, and appraised by
a verifier under a clear RATS-style role model~\cite{rats_rfc,eat_rfc,menetrey2022attestation}.

In RISC-V systems, this evidence chain is especially important because the
ecosystem encourages modular composition. A TVM may depend on a TSM, a platform
root of trust, an IOMMU, interrupt controllers, device-security protocols, and
firmware components from different vendors or projects. Attestation is the
mechanism that turns that composition into something a remote tenant can
evaluate.

\subsection{RISC-V Confidential I/O}

\subsubsection{Why CoVE Needs I/O Standardization}

A confidential VM that cannot use real devices is not sufficient for cloud or
edge deployment. CoVE-IO and TEE-I/O therefore extend the TVM discussion into
device identity, DMA, MMIO, interrupts, and link protection~\cite{riscv_cove_io_2026}. The
draft CoVE-IO specification is not ratified, but it is already useful as a map
of the required mechanisms. It introduces concepts such as trusted device
interfaces and device/security managers while relying on lower-level building
blocks including SPDM, TDISP, PCIe IDE, IOMMU support, and trusted interrupt
delivery~\cite{riscv_cove_io_2026,dmtf_spdm_2025,pcisig_tdisp_2022,pcie_ide}.
The TDISP citation is public PCI-SIG metadata for a member-gated ECN, so it
supports only the device-interface lifecycle role unless a stronger public
source is cited for a specific normative rule.

The design implication is straightforward: CoVE cannot stop at CPU-side page
ownership. A device assigned to a TVM must be identified, measured or certified,
placed into an appropriate trust state, constrained in the memory it can access,
and connected to the TVM through an interrupt path that does not silently return
control to the untrusted host. If any of those steps is missing, the I/O path
becomes a privileged channel around the confidential-computing boundary.

\subsubsection{RISC-V IOMMU, IOPMP, AIA, and sIOPMP}

The RISC-V IOMMU and AIA specifications provide two of the required supporting
layers. The IOMMU constrains device-side address translation and access
permissions, while AIA provides the interrupt and MSI architecture that CoVE-IO
can build on for trusted delivery~\cite{riscv_iommu_2023,riscv_aia_2023}. These
specifications are not confidential-computing systems by themselves. Their role
is to provide the architectural substrate on which a confidential I/O protocol
can express and enforce device ownership.

IOPMP and sIOPMP address a related but distinct problem: scalable filtering of
I/O transactions so that requester identity and memory-domain policy can be
checked outside the CPU's ordinary memory-access path~\cite{riscv_iopmp_2026,feng2024siopmp}.
For RISC-V confidential computing, this matters because DMA is the obvious way
to bypass a CPU-centered protection story. A TVM memory page protected from host
loads and stores is still exposed if a bus master can reach it without going
through a domain-aware check.

\subsection{Comparison with Arm CCA}

\begin{table}[htbp]
\centering
\caption{Arm CCA and RISC-V CoVE/AP-TEE Design Mapping}
\label{tab:arm_cca_riscv_cove}
\begin{tabular}{@{}p{0.27\linewidth}p{0.31\linewidth}p{0.31\linewidth}@{}}
\toprule
\textbf{Design Point} & \textbf{Arm CCA} & \textbf{RISC-V CoVE/AP-TEE} \\ \midrule
Protected workload & Realm & Trusted Virtual Machine \\
Management component & RMM & TSM and TSM driver \\
Host role & Resource management outside the Realm data boundary & Resource management outside the TVM data boundary \\
Memory ownership & GPT/GPC and Realm lifecycle & TVM memory donation, reclaim, sharing, and TSM-mediated state \\
Device path & SMMU and realm-aware device mediation & CoVE-IO, IOMMU, IOPMP, AIA, SPDM, TDISP, PCIe IDE \\
Standard status & Arm architecture/specification path & AP-TEE and CoVE-IO drafts, not ratified \\ \bottomrule
\end{tabular}
\end{table}

Table~\ref{tab:arm_cca_riscv_cove} shows why the comparison should be made at
the confidential-VM layer rather than between CCA and bare PMP. PMP remains a
useful primitive, especially for monitors and embedded TEEs, but CoVE/AP-TEE is
the closer architectural analogue to Arm CCA. Both systems try to let an
untrusted host manage resources without giving it data visibility. The
difference is that Arm presents a more vertically integrated Realm architecture,
whereas RISC-V standardization exposes a set of modular pieces that platform
builders must assemble into a coherent protection boundary.

\subsection{Summary}

RISC-V confidential computing has a rich pre-standardization lineage in
open-source enclave systems and a rapidly developing standardization path in
CoVE/AP-TEE. Earlier systems such as Sanctum, Keystone, TIMBER-V, CURE, MI6,
Penglai, SPEAR-V, Cerberus, and ACE establish the design vocabulary for
isolation, monitor structure, metadata, and memory sharing. CoVE/AP-TEE raises
the abstraction to confidential VMs, with TVM lifecycle, TSM mediation, memory
ownership transitions, and attestation as first-class concerns. CoVE-IO then
extends the same reasoning to devices and networks. The resulting lesson is
that RISC-V confidential computing is not one mechanism but a platform
composition problem: CPU privilege, memory ownership, I/O translation,
interrupt routing, device identity, and attestation all have to agree on the
same security boundary.


\section{Confidential I/O and Network Data-Path Defense}

Confidential computing protects data in use, but real workloads rarely remain
inside CPU registers and local DRAM. They use NICs, DPUs, SmartNICs, storage
controllers, GPUs, NPUs, CXL-attached memory, RDMA paths, and virtual switches.
These components move data through DMA engines, device queues, interrupts,
MMIO windows, link layers, and sometimes device-local runtimes. A CPU TEE or
confidential VM therefore provides only a partial boundary unless the I/O path
is integrated into the same trust model~\cite{riscv_cove_io_2026,arm_smmu_spec,sok-tee,zhu2020hetee,dhar2024cloudscale,zhou2024snic,giantsidi2025tnic}.

This section treats confidential I/O and network defense as a separate survey
lane. It does not cover generic network security. Instead, it focuses on
mechanisms that let a confidential workload safely use a device, network path,
remote memory tier, or accelerator without handing plaintext or control
authority back to an untrusted host.

\subsection{Threat Model and Data-Path Surfaces}

\subsubsection{Why Device Paths Break CPU-Only Isolation}

The basic device threat is simple: DMA-capable devices can access memory
without executing CPU load and store instructions. If the platform protects a
confidential VM only through CPU page tables, PMP, GPT/GPC, or monitor-managed
state, a device with insufficiently constrained DMA authority can still reach
protected pages~\cite{arm_smmu_spec,riscv_iommu_2023,riscv_iopmp_2026}. The
same concern applies to MMIO registers and queue metadata. A host that cannot
read a TVM's memory directly may still observe descriptors, packet buffers,
completion queues, or device status if those structures are not assigned and
protected as part of the confidential boundary.

Interrupts create a second surface. Device completion and error events must be
delivered to the protected workload without letting the host spoof, suppress, or
misroute security-critical signals. Work such as Devlore makes this explicit
for confidential VMs: interrupt ownership and delivery are part of device
ownership, not a secondary OS scheduling detail~\cite{bertschi2026devlore}.
For RISC-V, AIA and CoVE-IO provide the architectural vocabulary for trusted MSI
and interrupt handling, while Arm platforms rely on SMMU and interrupt
controller integration in the broader CCA system model~\cite{riscv_aia_2023,riscv_cove_io_2026,arm_smmu_spec}.

\subsubsection{Network and Offload-Specific Risks}

Network devices add their own complications. A confidential workload may use a
virtual NIC, SR-IOV virtual function, SmartNIC, DPU, programmable switch path,
or RDMA queue pair. In each case, the sensitive state is not only application
payload. It includes descriptors, keys, queue ownership, completion events,
flow steering rules, packet headers, and device firmware state. A secure
offload design must therefore answer four questions: which device or function
is the endpoint, what code and configuration is running there, which memory can
it access, and how packets or DMA buffers are protected while crossing the
fabric.

FOLIO illustrates one possible tradeoff for confidential VM networking by
aiming for high-performance networks without requiring trusted I/O
devices~\cite{li2024folio}. This is useful because it separates two design
paths. One path tries to make the device part of the trusted computing base
through assignment, measurement, and secure transport. The other path tries to
avoid trusting the device by keeping plaintext and sensitive control inside the
confidential VM while still obtaining high throughput. Both directions belong
in this survey because they address the same confidential-computing bottleneck:
network performance should not require silently expanding the trusted boundary
to an opaque device stack.

Network endpoint identity is a separate question from device identity. Even if
a NIC or DPU can be identified through SPDM-like mechanisms, a tenant still
needs to know whether the application channel carrying secrets terminates
inside the intended TEE software. TLS+RA addresses this endpoint problem by
binding remote attestation evidence to TLS channel establishment, so that a
secure channel is not merely adjacent to a TEE quote but cryptographically tied
to the attested endpoint~\cite{weinhold2025tlsra}. This matters for
confidential vNIC, proxy, gateway, and DPU-hosted services because TLS
termination can otherwise become a place where plaintext exits the confidential
boundary.

SmartNIC and NIC-local trust work pushes the same question lower in the stack.
\textbf{S-NIC} shows that a multi-tenant SmartNIC needs hardware isolation for
on-NIC RAM, cache, accelerators, DMA, and bus bandwidth before offloaded network
functions can be treated as tenant-confidential components~\cite{zhou2024snic}.
\textbf{TNIC} takes a different route by placing a minimal root of trust at the
NIC, providing transferable authentication and non-equivocation for trustworthy
distributed systems~\cite{giantsidi2025tnic}. These works are not substitutes
for VM memory protection, SPDM, or TDISP. They clarify what a NIC or DPU must
provide if it becomes a real endpoint of a confidential data path.

\subsection{Protocol Layer: Identity, Session Security, and Device State}

\subsubsection{SPDM as Device Identity and Measurement Substrate}

The \textbf{Security Protocol and Data Model (SPDM)} provides a protocol for
device authentication, measurement retrieval, version and capability
negotiation, and session establishment~\cite{dmtf_spdm_2025,dmtf_spdm_arch_2024}.
In confidential-computing systems, SPDM is valuable because the host's view of a
device is not enough. A tenant or trusted manager needs evidence about the
device identity and state before allowing that device to participate in a
protected data path.

SPDM secured messages add confidentiality and integrity for protocol exchanges
after session establishment~\cite{dmtf_spdm_secured_messages_2025}. SPDM over
TCP extends the transport setting so that SPDM can be used beyond only local
bus-oriented interactions~\cite{dmtf_spdm_tcp_2024}. These specifications do
not themselves define a complete confidential I/O architecture. Their role is
to supply the identity, measurement, and secure-channel substrate that higher
level mechanisms such as TDISP, CoVE-IO, and trusted I/O designs can rely on.

\subsubsection{TDISP and Trusted Device Interfaces}

\textbf{TDISP} addresses a more specific lifecycle problem: how a PCIe device
interface can be placed into a TEE-aware trust state and assigned to a
confidential workload~\cite{pcisig_tdisp_2022}. Device identity alone is not
enough. The platform also needs a way to bind a particular device interface,
function, or trusted device interface to the workload that will use it. TDISP
therefore belongs at the device-interface lifecycle layer, whereas SPDM belongs
at the identity and secure-session layer.

The RISC-V CoVE-IO draft uses this kind of structure explicitly. It refers to
trusted device interfaces, device managers, security managers, SPDM, TDISP,
PCIe IDE, IOMMU controls, and trusted MSI as cooperating parts of a
confidential I/O design~\cite{riscv_cove_io_2026}. The important survey point
is that no single layer is sufficient. A measured device without constrained
DMA can still leak memory. A constrained DMA path without device identity can
still attach the wrong endpoint. A protected link without lifecycle control can
still carry traffic for an interface that was not securely assigned.

\subsubsection{Attested Service Channels}

SPDM and TDISP focus on the device and device-interface layers, but confidential
workloads also expose application-level endpoints. A key broker, model-serving
endpoint, storage proxy, packet-processing service, or DPU-hosted TLS
termination service may all need to prove that the network channel reaches the
expected protected code. \textbf{TLS+RA} integrates remote attestation into the
TLS handshake while keeping certificate-based endpoint identity and
attestation-based platform evidence independently deployable and independently
failing~\cite{weinhold2025tlsra}. This gives the survey a seventh layer:
attested service-channel binding. It is distinct from device identity because
the question is not only "which device is present" but "where does this
application connection terminate."

\subsection{Data Plane: DMA, Links, and Fabrics}

\subsubsection{IOMMU, SMMU, IOPMP, and Address Authority}

DMA mediation is the first enforcement layer for confidential I/O. Arm systems
use SMMU-based translation and protection to constrain device accesses in the
system memory path~\cite{arm_smmu_spec}. RISC-V systems use the IOMMU
specification for device address translation and IOPMP-style mechanisms for
requester-aware transaction filtering~\cite{riscv_iommu_2023,riscv_iopmp_2026}.
Academic work such as sIOPMP emphasizes the scalability pressure: protected
systems may have many devices, many domains, and frequent mapping changes, so
the I/O protection mechanism must be both precise and efficient~\cite{feng2024siopmp}.

These mechanisms should be classified as access-control and translation
mechanisms. They do not automatically encrypt data on a link, validate the
device's internal firmware, or prove the endpoint state to a remote verifier.
Their job is to ensure that a bus master can access only the memory and MMIO
resources assigned to it. This classification prevents a common survey error:
equating any I/O protection with a complete trusted I/O system.

\subsubsection{PCIe IDE, CXL, RDMA, and Remote Memory Paths}

Link and fabric protection answer a different question: what happens to data as
it traverses PCIe, CXL, RDMA, or related interconnects. PCIe IDE provides
integrity and data encryption for PCIe traffic~\cite{pcie_ide}. CXL changes the
memory hierarchy by allowing memory expansion, pooling, and tiering beyond
local DRAM~\cite{cxl_spec,gouk2022directcxl,zhong2024cxltiers}. RDMA and
remote paging systems move application memory over the network or fabric with
minimal CPU involvement~\cite{wang2025odrp}. These technologies are not all
confidential-computing defenses by themselves, but they define the paths that a
confidential-computing defense must cover.

The design consequence is that confidential workloads need a data-path view, not
only a CPU view. A page that is protected while resident in local memory may
later move through a CXL fabric, be staged through an RDMA-capable NIC, or be
referenced by a device queue. Each transition introduces a new enforcement
point. The platform must decide whether protection comes from link encryption,
DMA access control, endpoint TEE support, software copying through trusted
buffers, or a combination of these mechanisms.

Secure storage over network fabrics adds freshness and crash-consistency
questions to this path. \textbf{Hazel} studies NVMe-over-Fabrics storage under a
confidential-computing threat model and combines counter leasing, NVMe metadata,
a disaggregated Merkle-tree design, and BlueField-3 offload to provide
confidentiality, integrity, and freshness without changing the NVMe-oF
protocol~\cite{chrapek2026hazel}. This kind of work is useful for the survey
because it shows that a protected I/O path must sometimes defend not only
in-flight packets, but also remote persistent state and replay freshness.

\subsection{Trusted I/O Systems and Device TEEs}

\subsubsection{SEV-TIO and VM-Oriented Trusted I/O}

AMD SEV-TIO provides a concrete industry design point for trusted I/O in
encrypted-VM systems~\cite{amd_sev_tio_2023}. It is useful in this survey even
though the main architectural focus is Arm and RISC-V, because it clarifies the
general pattern: a confidential VM needs device assignment, device trust
establishment, protected DMA, interrupt handling, and attestation-compatible
state. These requirements reappear in CoVE-IO and in Arm CCA-oriented device
work.

The SEV-TIO comparison also helps avoid an overly narrow interpretation of
trusted I/O. Trusted I/O is not merely "attach a device to a VM." It is the
combination of device identity, interface lifecycle, memory access restriction,
link security, interrupt delivery, and evidence that a verifier can evaluate.
Different platforms divide those responsibilities differently, but the required
security properties are similar.

\subsubsection{Accelerator TEEs and Confidential Offload}

Accelerator TEEs push the same problem into device-local computation. The
accelerator may execute kernels, hold model weights, stage tensors, manage
queues, or perform packet-processing offload. A CPU confidential VM cannot
automatically protect that state once it leaves the CPU package. The accelerator
itself must either become part of the trusted boundary or be treated as
untrusted with protocols that keep secrets inside the confidential workload.

ACAI studies accelerator protection in the Arm CCA context~\cite{acai2023}.
PORTAL studies protected Arm CCA device access for mobile SoCs, using CCA
memory isolation to avoid repeated shared-buffer encryption~\cite{sang2025portal}.
CAGE adapts GPU and FPGA accelerator workflows to Arm CCA with shadow tasks and
GPC/GPT-backed checks~\cite{wang2026cage}. StrongBox is an earlier Arm GPU TEE
baseline that uses TrustZone and Stage-2 translation to protect unified-memory
GPU computation~\cite{deng2022strongbox}.
ITX demonstrates confidential computing within an AI accelerator, giving a
concrete example of device-local confidential execution~\cite{vaswani2023itx}.
The accelerator TEE SoK organizes this space around device identity, memory and
bus protection, runtime TCB, scheduling, and attestation~\cite{sok-tee}. These
works are important for network defense as well because modern NICs, DPUs, and
SmartNICs increasingly resemble accelerators: they execute firmware and
programmable code, maintain queues, and transform data before the CPU sees it.

\subsubsection{SmartNIC, DPU, and Rack-Scale Offload}

SmartNICs and DPUs sit between the device-interface layer and the application
service layer. They can expose SR-IOV functions and virtual switches to the
host, run embedded Arm-side control software, terminate tunnels, process
packets, or mediate storage traffic. HETEE and later cloud-scale work show one
route for protecting accelerator and DPU-like offload: place a security
controller on the data path so that non-TEE devices can be assigned, measured,
and cleaned up as part of a confidential workload lifecycle~\cite{zhu2020hetee,dhar2024cloudscale}.
Commercial DPU documentation shows another route: DPU-local roots such as
OP-TEE/fTPM and replay-protected storage can provide platform identity and key
protection building blocks~\cite{nvidia_bluefield_operation_2025}. That
commercial evidence should be used carefully, because the cited BlueField
documentation describes fTPM support and its limitations rather than a complete
production confidential networking TEE.

S-NIC and TNIC fill two additional gaps. S-NIC answers how a SmartNIC can
isolate tenant network functions from other functions and from the NIC OS, while
TNIC answers how a NIC can provide a small, verifiable trust root for networked
protocols~\cite{zhou2024snic,giantsidi2025tnic}. Hazel then shows how DPU/SmartNIC
offload can be used in a secure storage data path while preserving freshness
properties~\cite{chrapek2026hazel}. Together these papers make the DPU/NIC
category more precise: one must distinguish local resource isolation, endpoint
attestation, protocol-level non-equivocation, and storage freshness.

For survey classification, these systems are neither ordinary firewall/IDS
papers nor complete replacements for CoVE-IO or TDISP. They are evidence that
network offload devices are becoming programmable trust endpoints. A secure
confidential networking stack must therefore compose DPU-local roots,
device-interface lifecycle, DMA restrictions, link protection, and attested
application-channel termination.

\subsection{Taxonomy}

\begin{table*}[t]
\centering
\caption{Confidential I/O and Network Defense Taxonomy}
\label{tab:confidential_io_taxonomy}
\begin{tabular}{@{}p{0.18\linewidth}p{0.26\linewidth}p{0.25\linewidth}p{0.22\linewidth}@{}}
\toprule
\textbf{Layer} & \textbf{Primary Question} & \textbf{Representative Mechanisms} & \textbf{Survey Role} \\ \midrule
Device identity and measurement & Which endpoint is participating, and what state can it prove? & SPDM, SPDM secured messages, SPDM over TCP & Establish evidence and secure sessions for device trust. \\
Device-interface lifecycle & Which interface is assigned to the confidential workload? & TDISP, trusted device interface state, CoVE-IO TDI/TDM/DSM & Bind a device function or interface to a protected domain. \\
Attested service channel & Does the application channel terminate in the attested protected endpoint? & TLS+RA, RA-TLS-style endpoint binding, verifier policy & Bind network services, KBS endpoints, proxies, and DPU-hosted TLS termination to TEE evidence. \\
DMA and MMIO authority & Which memory and registers may the device access? & SMMU, RISC-V IOMMU, IOPMP, sIOPMP & Enforce address translation and transaction filtering. \\
Interrupt and event delivery & Can completion and error events be delivered without host spoofing or confusion? & AIA, trusted MSI, Devlore-style interrupt ownership & Keep device ownership coherent after asynchronous events. \\
Link and fabric protection & What protects traffic across PCIe, CXL, RDMA, and remote-memory paths? & PCIe IDE, CXL fabric controls, protected RDMA/remote paging designs & Protect data once it leaves CPU-local memory paths. \\
NIC-local root and function isolation & Can a SmartNIC/DPU isolate tenant functions or provide a trusted network primitive? & S-NIC, TNIC, DPU OP-TEE/fTPM & Treat NICs and DPUs as explicit trust endpoints rather than opaque offload boxes. \\
Confidential storage data path & Can remote storage preserve confidentiality, integrity, and freshness? & Hazel, protected NVMe-oF metadata, DPU crypto offload & Extend confidential I/O beyond packets into persistent remote state. \\
Device-local confidential execution & Does the accelerator or offload engine hold secrets or execute trusted code? & ACAI, PORTAL, CAGE, StrongBox, ITX, HETEE, CloudScale, BlueField OP-TEE/fTPM, accelerator TEE designs & Extend or avoid extending the TCB to accelerators and offload devices. \\ \bottomrule
\end{tabular}
\end{table*}

Table~\ref{tab:confidential_io_taxonomy} captures the classification used in
the rest of this survey. The table is intentionally layered because
confidential I/O failures are often composition failures. A system can have
SPDM but weak DMA protection, PCIe IDE but no trusted device-interface
lifecycle, or a device TEE but an ambiguous interrupt path. The practical
question is whether all relevant layers preserve the same ownership boundary.

\subsection{Summary}

Confidential I/O and network defense is the point where confidential computing
becomes a whole-platform problem. SPDM and secured messages provide identity,
measurement, and session foundations. TDISP and CoVE-IO introduce trusted
device-interface lifecycle concepts. TLS+RA-style protocols bind application
channels to attested endpoints. SMMU, IOMMU, IOPMP, and sIOPMP constrain DMA
authority. AIA and trusted MSI address interrupt delivery. PCIe IDE, CXL, RDMA,
NVMe-oF, and remote-memory systems define the fabric paths that protected data
may cross. Accelerator, SmartNIC, NIC, and DPU work such as ACAI, PORTAL, CAGE,
StrongBox, ITX, HETEE, CloudScale, S-NIC, TNIC, Hazel, and BlueField
OP-TEE/fTPM shows that device-local, NIC-local, storage-path, or
controller-mediated execution can become part of the confidential boundary. A
survey that focuses only on CPU TEE mechanisms would miss this entire class of
defenses; for modern Arm and RISC-V confidential-computing platforms, these
mechanisms are core rather than peripheral.


\section{Security Design Mechanisms in Arm and RISC-V}

Hardware security mechanisms are increasingly being exposed as architectural primitives rather than as optional software conventions. Across modern Arm and RISC-V systems, security-relevant invariants such as indirect-control validation, page-table authority, physical-memory ownership, and privileged-access constraints are moving into the hardware-software contract itself~\cite{armv-a,riscv_privileged}. For a survey of hardware security, this shift matters because it changes where defenses live: instead of relying only on compiler instrumentation or operating-system policy, the architecture now provides direct mechanisms for constraining control flow, limiting memory authority, and mediating access across privilege and device boundaries. The following sections survey these mechanisms by design objective, focusing on runtime integrity, privilege and memory isolation, memory-safety hardening, and the system-integration issues that determine whether a mechanism remains effective once devices, debug paths, and accelerators are considered.

\subsection{Control-Flow and Memory-Safety Mechanisms}

\subsubsection{Arm Runtime-Integrity Mechanisms}

Arm's recent A-profile architecture incorporates several mechanisms that aim to make common exploitation primitives harder to realize in practice~\cite{armv-a}. \textbf{Branch Target Identification (BTI)} constrains indirect branches to architecturally valid landing pads, reducing the usable gadget space for jump-oriented and related control-flow attacks. \textbf{Pointer Authentication (PAC)} protects pointers with cryptographic authentication codes, making pointer substitution and reuse attacks more difficult. \textbf{Guarded Control Stack (GCS)} adds hardware support for protecting return-address state. These mechanisms address different parts of the control-flow problem, but together they represent a clear design direction: indirect control transfer is no longer treated as something that software may validate only if it chooses to.

Arm applies a similar idea to memory-safety failures. \textbf{Memory Tagging Extension (MTE)} associates tags with memory granules and checks tag consistency on pointer dereference, allowing the hardware to detect classes of spatial and temporal memory-safety violations with relatively low overhead~\cite{armv-a}. From a survey perspective, the important point is not that any one mechanism eliminates memory corruption, but that the architecture increasingly offers built-in support for narrowing the gap between a latent software bug and a reliable control-hijacking exploit.

\subsubsection{RISC-V Control-Flow and Protection Primitives}

The RISC-V architectural material surveyed here shows a related, though more modular, direction. The privileged and associated architectural specifications include support for control-flow integrity through extensions such as \textbf{Zicfilp}, which constrains indirect-control targets through landing-pad semantics, and \textbf{Zicfiss}, which provides hardware support for protecting return-address state~\cite{riscv_privileged}. Rather than embedding one vertically integrated protection model, the RISC-V ecosystem exposes a set of composable primitives that can be incorporated into different classes of systems.

This modularity is particularly visible in the treatment of memory authority. In the cited RISC-V specifications, the most established protection story centers on privilege separation and physical-memory access control, especially through \textbf{Physical Memory Protection (PMP)} and its extensions, rather than on a single Arm-MTE-style mainstream memory-tagging mechanism~\cite{riscv_privileged}. The result is a design space in which control-flow hardening and memory isolation exist, but are often assembled as part of a larger platform composition.

\subsection{Privilege, Physical Memory, and Domain Isolation}

\subsubsection{Arm Privilege Hierarchy and System-Wide Isolation}

Arm's protection model combines conventional privilege separation with system-wide domain isolation. The exception-level hierarchy from EL0 to EL3 separates application, kernel, hypervisor, and monitor software, while \textbf{TrustZone} extends this structure with Secure and Non-secure worlds enforced across the SoC~\cite{arm_trustzone_whitepaper,pinto2019trustzone,armv-a}. Additional mechanisms such as \textbf{PAN} and \textbf{PXN} constrain privileged software inside a single world by preventing inappropriate access to or execution from lower-privilege mappings~\cite{armv-a}. These mechanisms are important because many real attacks do not begin by crossing a world boundary; they begin by exploiting excessive authority inside an already privileged context.

At the system level, Arm's more recent confidential-computing extensions push isolation further into the memory system. \textbf{Arm CCA} and \textbf{RME} use hardware checks over physical-memory ownership to protect \textbf{Realms} even in the presence of a powerful host stack~\cite{li2022cca,arm_cca_spec,arm_rme_spec}. In other words, the architecture is moving from privilege-only isolation toward coordinated CPU, memory, and device mediation. This progression is part of a broader trend in secure hardware design: once DMA-capable devices and privileged software are part of the threat model, core-local access checks are not enough.

\subsubsection{RISC-V Privilege Modes and PMP-Based Isolation}

RISC-V organizes protection around privilege modes and a register-programmed physical-memory filter. \textbf{PMP} allows machine mode to define access permissions for physical-memory regions, thereby constraining S-mode and U-mode software~\cite{riscv_privileged}. Extensions such as \textbf{ePMP} and \textbf{Smepmp} strengthen this model by refining the handling of execute permissions and reducing the amount of implicitly trusted machine-mode behavior~\cite{riscv_privileged}. The design is comparatively simple and well suited to embedded systems, monitors, and smaller TEEs where a limited number of protected regions is acceptable.

The main limitation of this style appears when system complexity increases. Large SoCs, confidential-computing platforms, and accelerator-rich systems must often coordinate CPU isolation with \textbf{DMA} control and device-visible memory ownership. Recent work such as \textbf{sIOPMP} therefore studies how scalable I/O protection can complement CPU-side isolation for TEEs~\cite{feng2024siopmp}. This is an important design point for survey work: in RISC-V platforms, the strength of the final security story depends heavily on how PMP-like CPU mechanisms are composed with peripheral protection.

\subsection{Page Tables, Capabilities, and Memory-Protection Semantics}

\subsubsection{PTE and Page-Table Permissions}

Page tables are a first-order hardware defense mechanism, not only a virtual-memory convenience. In Arm systems, page-table attributes and privileged-access controls define whether a mapping is readable, writable, executable, accessible from lower privilege, or usable as an instruction source~\cite{armv-a}. Mechanisms such as \textbf{PXN}, \textbf{UXN}, and \textbf{PAN} are therefore best classified as page-table and privilege-hardening defenses. They reduce the ambient authority of privileged software by making kernel execution from user-controlled mappings and privileged access to user memory explicit policy choices rather than default behaviors.

RISC-V has an analogous but differently organized permission story. RISC-V page-table entries encode read, write, execute, user, global, accessed, and dirty state, while PMP/ePMP/Smepmp provide physical-address constraints below or beside the virtual-memory layer~\cite{riscv_privileged}. The important survey distinction is that page-table permissions, PMP-style filters, and confidential-computing ownership metadata are different layers. Page tables express address-translation and per-page software protection policy. PMP and IOPMP constrain physical-memory and transaction authority. CCA GPT/GPC and CoVE/AP-TEE lifecycle state express confidential-domain ownership. Confusing these layers leads to incorrect claims about what a mechanism actually defends.

\subsubsection{Capability and Tagging Architectures}

Capability and tagging architectures take a more semantic approach to memory authority. \textbf{CHERIoT} applies capability-based memory safety to embedded devices and demonstrates how bounds, permissions, and object provenance can be represented in hardware-enforced capabilities~\cite{amar2023cheriot}. \textbf{RV-CURE} explores a RISC-V capability architecture for full memory safety~\cite{kim2023rvcure}. These mechanisms are relevant to confidential-computing platforms because many attacks on protected workloads still begin as ordinary memory-safety failures inside the trusted or protected domain. A Realm or TVM can hide data from the host while still being vulnerable to corruption inside its own software stack.

The correct classification is therefore defense-in-depth. Capability and tagging mechanisms do not replace CCA, CoVE, TrustZone, or device assignment. Instead, they reduce the probability that code inside a protected domain can be exploited to misuse the domain's legitimate authority. This is a recurring theme in modern hardware security: isolation defines who is outside the boundary, while memory-safety and capability mechanisms reduce abuse by code that is already inside it.

\subsubsection{Memory Encryption, Integrity, and Replay Protection}

Memory encryption should also be separated from access-control mechanisms. A memory-encryption design addresses what an observer of raw memory or a memory bus can learn, while integrity and replay-protection designs address whether modified or stale data can be detected~\cite{henson2014memory}. AMD SEV-SNP is an important modern example because it combines encrypted-VM deployment with stronger protection against host-controlled memory-management attacks~\cite{amd_sev_snp}. Arm CCA and RISC-V CoVE/AP-TEE, by contrast, primarily specify ownership, lifecycle, and access-control semantics in the cited architectural materials~\cite{arm_cca_spec,riscv_ap_tee_2024}. A complete product may combine these layers, but a survey should not treat them as interchangeable.

This taxonomy is especially important for Arm and RISC-V confidential-computing comparison. GPT/GPC, PMP, IOPMP, and IOMMU checks decide whether an access is allowed. Link protection such as PCIe IDE protects traffic on a fabric. Memory encryption protects data representation outside the processor or protected boundary. Integrity and anti-replay mechanisms protect freshness and tamper detection. These mechanisms compose, but none of them alone proves that a workload is confidential against the full host, device, and fabric threat model.

\subsection{Trusted Execution, Boot, and Trust Anchors}

\subsubsection{TEE-Oriented Design Patterns}

Trusted-execution support remains one of the clearest ways in which hardware security mechanisms appear as complete platform designs rather than isolated instructions. In the Arm ecosystem, \textbf{TrustZone} provides the classical two-world TEE model and has been used as the basis for a large body of systems work and security analysis~\cite{arm_trustzone_whitepaper,pinto2019trustzone,guan2017trustshadow,cerdeira2020trustzone}. Arm CCA extends this line further by introducing \textbf{Realms} for confidential workloads under stronger host-side adversary assumptions~\cite{li2022cca,arm_cca_spec,arm_rme_spec}. More broadly, the comparison between TrustZone-style TEEs and CCA-style confidential execution shows how trusted-execution design has moved from isolating a privileged security service toward protecting workloads against stronger software adversaries.

\subsubsection{Boot-Time Trust and Measured Launch}

Trusted execution is closely connected to how a platform establishes trust at boot time. Work on secure boot, trusted boot, and remote attestation for Arm TrustZone-based devices demonstrates that the security of the runtime environment depends on a measured launch chain rather than on isolation alone~\cite{ling2021trustzoneatt,chen2024mraima,mao2025pdrima}. At a broader systems level, \textbf{SWATT}, \textbf{SMART}, and \textbf{SEDA} show recurring design requirements for attestation and roots of trust: measurements must be attributable, fresh, and tied to a concrete execution context~\cite{seshadri2004swatt,defrawy2012smart,asokan2015seda}. For hardware designers, this means that the relevant trust anchor is not simply a privileged mode or a secure memory region; it is the combination of launch-time measurement, protected execution, and verifiable system state.

\subsection{I/O, DMA, and Accelerator Mediation}

\subsubsection{Arm System Mediation Beyond the CPU}

Once devices and accelerators are part of the threat model, hardware security depends on more than CPU privilege checks. Arm's system architecture therefore includes I/O mediation through the \textbf{SMMU}, which is especially important when confidential workloads or secure services interact with \textbf{DMA}-capable peripherals~\cite{arm_smmu_spec}. The same point is visible in Arm CCA-oriented systems work, where accelerator protection and device-interrupt delivery are treated as part of the core isolation problem rather than as optional external policy~\cite{acai2023,bertschi2026devlore}. This is a crucial design lesson for secure hardware surveys: a platform that protects the CPU but leaves devices outside the security model has only partial isolation.

\subsubsection{RISC-V I/O Protection as a Complementary Layer}

RISC-V reaches the same requirement through a more modular route. The base privileged architecture centers on CPU-side protection, while later work such as \textbf{sIOPMP} studies scalable I/O protection as a complementary mechanism for TEEs and protected domains~\cite{riscv_privileged,feng2024siopmp}. This separation reinforces the broader pattern seen throughout the chapter: RISC-V often exposes a toolkit whose final security properties depend on disciplined composition, whereas Arm more often presents these pieces as parts of a vertically integrated platform-security story.

\subsection{Attestation and Security Lifecycle Mechanisms}

\subsubsection{Attestation as a Hardware-Supported Design Mechanism}

Attestation is not only a protocol feature; it is also a hardware-design concern because the architecture must expose stable roots of trust, measurable state, and protection domains whose claims can be verified. The \textbf{EAT} standard provides a general structure for carrying attestation claims, while broader work on attestation mechanisms for TEEs emphasizes that the usefulness of those claims depends on what the hardware and trusted runtime are able to measure and report~\cite{eat_rfc,menetrey2022attestation}. In this sense, attestation belongs in a survey of hardware security design because it links platform mechanisms to external trust decisions.

\subsubsection{Lifecycle Security Frameworks}

The same design logic appears in platform frameworks such as \textbf{PSA Certified}, which treats hardware-rooted security, attestation, and lifecycle management as part of one assurance story rather than as isolated controls~\cite{psa_certified}. This kind of framework is important because it extends the discussion beyond the presence of security mechanisms to their use across provisioning, update, and operational phases. For surveys of secure hardware design, lifecycle treatment matters because many platforms fail not through the absence of a mechanism, but through inconsistent handling of that mechanism across deployment stages.

\subsection{Debug, Trace, and System Integration}

\subsubsection{Security-Relevant Auxiliary Components}

Architectural security mechanisms are not limited to branches, page tables, and privilege bits. Arm's documentation on trace makes clear that debug and observability features such as \textbf{ETM} provide detailed execution visibility that is valuable for validation and diagnostics, but that same visibility must be tightly controlled in security-sensitive deployments~\cite{arm_understanding_trace}. This illustrates a broader systems lesson: every auxiliary subsystem that can observe execution or preserve state across contexts can become part of the attack surface if it is not incorporated into the protection model.

\subsubsection{Composition as the Central Design Problem}

Across both Arm and RISC-V, the mechanisms surveyed above point to the same conclusion. Secure hardware design is not achieved by enabling a single extension in isolation. Control-flow hardening, memory-access constraints, privilege separation, and device mediation have to compose correctly across the whole platform~\cite{armv-a,arm_cca_spec,arm_rme_spec,riscv_privileged,feng2024siopmp}. This is also why architecture manuals and systems papers increasingly discuss memory ownership, I/O paths, and monitor logic in the same breath: the boundary that matters in practice is the one that survives contact with the full machine, not just the CPU pipeline.

\subsection{Cross-Architecture Comparison}

Table~\ref{tab:cross_arch_security_mapping} summarizes the main mechanism categories surveyed in this section. The goal is not to claim that the two ecosystems are identical, but to show that they increasingly address comparable security objectives through different integration styles.

\begin{table}[htbp]
\centering
\caption{Cross-Architecture Security Mechanism Mapping}
\label{tab:cross_arch_security_mapping}
\begin{tabular}{@{}p{0.28\linewidth}p{0.30\linewidth}p{0.30\linewidth}@{}}
\toprule
\textbf{Mechanism Category} & \textbf{Arm} & \textbf{RISC-V} \\ \midrule
Forward-edge control-flow protection & BTI & Zicfilp \\
Return-address protection & GCS / PAC-assisted hardening & Zicfiss \\
Memory bug detection & MTE & Protection centered on privilege and memory-access control in the cited materials \\
Primary privilege structure & EL0--EL3 plus Secure/Non-secure system partitioning & M/S/U privilege modes \\
Physical-memory isolation & TrustZone system partitioning; CCA/RME ownership checks & PMP, ePMP, Smepmp \\
Page-table hardening & PAN, PXN, UXN, privileged translation controls & PTE R/W/X/U/A/D permissions and privilege-aware translation \\
Capability or tagging memory safety & MTE and capability-adjacent ecosystem work & CHERIoT, RV-CURE, tag/capability research systems \\
System-level I/O protection & SMMU-mediated system integration in Arm platforms & IOMMU, IOPMP, sIOPMP, AIA-based interrupt support \\ \bottomrule
\end{tabular}
\end{table}

The comparison suggests a consistent difference in emphasis. The Arm material surveyed here presents a more vertically integrated security stack in which runtime integrity, world/domain isolation, and system-wide mediation are documented as parts of one platform story~\cite{armv-a,arm_cca_spec,arm_rme_spec}. The RISC-V material, by contrast, emphasizes composable building blocks centered on privilege modes and PMP-based physical-memory control, with additional system protections layered on through complementary mechanisms and proposals~\cite{riscv_privileged,feng2024siopmp}. This is not a value judgment so much as a design distinction: one ecosystem foregrounds integration, while the other foregrounds modular composition.

\subsection{Implications for Secure Hardware Design}

For hardware designers, the practical lesson is that secure systems must be evaluated as whole-machine constructions. A platform with strong control-flow protections can still fail if debug visibility is unmanaged. A platform with strict CPU isolation can still fail if a device bypasses memory protections through \textbf{DMA}. A platform with clean privilege layering can still fail if the architecture leaves too much ambient authority in the most privileged software context. The mechanisms surveyed in this section therefore support the same conclusion reached elsewhere in this paper: modern hardware security depends on coordinated enforcement across the instruction set, the memory system, and the device path.

\subsection{Summary}

The architectures surveyed here show a broad movement toward hardware-enforced security invariants. In Arm, this is visible in the combination of BTI, PAC, GCS, MTE, page-table hardening, privilege constraints, TrustZone, and CCA/RME-based system isolation~\cite{armv-a,arm_trustzone_whitepaper,arm_cca_spec,arm_rme_spec}. In RISC-V, the same trend appears through control-flow integrity extensions, PMP-family physical-memory protection, page-table permissions, capability/tagging research, and complementary proposals for scalable I/O mediation~\cite{riscv_privileged,feng2024siopmp,kim2023rvcure,amar2023cheriot}. A cross-cutting observation is that the core design problem is no longer whether hardware offers security features at all, but whether those features compose into a coherent platform-level protection model once privileged software, accelerators, and device traffic are part of the system.


\bibliographystyle{IEEEtran}
\bibliography{reference}

\end{document}
